\documentclass[11pt]{article}
\usepackage[margin=1in]{geometry}
\usepackage{longtable}
\usepackage{booktabs}
\usepackage{hyperref}
\usepackage{xcolor}
\usepackage{amsmath}

\title{Independent Replication Report --- BVBRC-72\\
\textit{Bacillus amyloliquefaciens} WS-8: Genomics and Secondary Metabolite BGC Confirmation}
\author{OpenClaw Replication Pipeline}
\date{2026-07-06}

\begin{document}
\maketitle

\begin{abstract}
This report is the independent-replication analysis of Liu et al.\ 2020 (\textit{J.\ Microbiol.\ Biotechnol.}\ 30(3):417--426, DOI \href{https://doi.org/10.4014/jmb.1906.06055}{10.4014/jmb.1906.06055}, PMID 31601062, PMC 9728402), which reported the PacBio SMRT genome assembly, PGAP annotation, antiSMASH v3.0 BGC prediction, LC-ESI-Q-TOF-MS lipopeptide profiling, and late-log RNA-Seq of the plant-growth-promoting rhizobacterium \textit{Bacillus amyloliquefaciens} strain WS-8 (GenBank CP018200.1 / GCF\_001922005.1). Verdict: \textbf{PARTIAL $\rightarrow$ REPLICATED-leaning} (LLM-judge coverage 90\,\%, agreement 83.6\,\%). All PGAP-derived quantitative genome claims and all seven named biosynthetic gene clusters reproduced exactly using the deposited assembly + fresh antiSMASH v8.0.4 KnownClusterBlast; the two wet-lab-only claim families (LC-MS metabolomics, RNA-Seq expression) could not be re-executed because the raw datasets were never deposited.
\end{abstract}

\section{Paper summary}
Liu et al.\ PacBio-sequenced \textit{B.\ amyloliquefaciens} WS-8 to a single circular chromosome of 3{,}929{,}787\,bp with no plasmid, annotated it via NCBI PGAP, and ran antiSMASH 3.0 to predict 19 secondary-metabolite biosynthetic gene clusters (BGCs). Seven of those clusters showed $>$70\,\% gene similarity to reference BGCs for the antifungal lipopeptides/polyketides difficidin, fengycin, bacillaene, macrolactin, surfactin, bacilysin, and the siderophore bacillibactin, plus one class-II lanthipeptide cluster framed as novel. LC-ESI-Q-TOF-MS on WS-8 culture extracts identified 21 lipopeptide species (5 iturins, 16 fengycins) as the dominant antifungal principle against \textit{Botrytis cinerea}.

\section{Claims and outcomes}

\begin{longtable}{p{0.7cm}p{6.7cm}p{1.6cm}p{1.5cm}p{4.5cm}}
\toprule
\textbf{\#} & \textbf{Claim} & \textbf{Type} & \textbf{Result} & \textbf{Metric} \\
\midrule
\endhead
C1 & Genome = 3{,}929{,}787\,bp, single circular chromosome & quant & \checkmark & 3{,}929{,}787\,bp; circular \\
C2 & No plasmid & qual & \checkmark & 1 replicon \\
C3 & GC = 46.5\,\% (abstract) / 45.6\,\% (Table 1) & quant & \checkmark abstract; Table 1 typo & Direct: 46.499\,\% \\
C4 & 3895 predicted genes & quant & \checkmark & 3895 \\
C5 & 3777 CDS & quant & \checkmark & 3777 \\
C6 & 107 pseudogenes & quant & \checkmark & 107 \\
C7 & 86 tRNA & quant & \checkmark & 86 \\
C8 & 27 rRNA (9 operons 5S/16S/23S) & quant & \checkmark & 9$\times$5S+9$\times$16S+9$\times$23S \\
C9 & 5 ncRNA & quant & $\sim$\checkmark & 4 ncRNA + 1 tmRNA \\
C10 & 19 BGCs (antiSMASH 3.0) & quant & tool-delta & 13 (v8) / 12 (v7) \\
C11 & 1 class-II lanthipeptide BGC & qual & \checkmark & v8 Region 8; no MIBiG hit \\
C12 & Fengycin + bacillaene $>$70\,\% & qual & \checkmark & 97--99\,\% identity, both \\
C13 & Macrolactin $>$70\,\% & qual & \checkmark & 97--99\,\% identity \\
C14 & Difficidin $>$70\,\% & qual & \checkmark & 97--99\,\% identity \\
C15 & Bacilysin $>$70\,\% & qual & \checkmark & 98--100\,\% identity \\
C16 & Bacillibactin $>$70\,\% & qual & \checkmark(variant) & 57--81\,\% (dhbA-F variant) \\
C17 & Surfactin $>$70\,\% & qual & \checkmark & 95--100\,\% identity \\
C18 & LC-MS: iturins + fengycins dominate & wet-lab & spot-check & Iturin BGC \emph{also} present \\
C19 & RNA-Seq: 6 antifungal BGCs expressed & wet-lab & spot-check & No SRA; genes present \\
C20 & Auxin biosynthesis genes & qual & $\sim$\checkmark & Trp operon + aromatic AAT \\
C21 & PacBio SMRT, $\sim$311$\times$ coverage & meta & \checkmark platform; unverifiable depth & SRA not deposited \\
\bottomrule
\end{longtable}

\noindent Coverage: 21/21 = 100\,\%; agreement 83.6\,\% (LLM-judge \texttt{argo:gpt-5.2}).

\section{Method}
Full method is at \S3 of \texttt{REPORT.md}. Summary:
\begin{enumerate}
\item Retrieved paper via Europe PMC (PMC9728402) + PubMed efetch (PMID 31601062).
\item Identified assembly GCF\_001922005.1 (CP018200.1) via NCBI Assembly; exact bp-length match to paper.
\item Downloaded GenBank + FASTA via NCBI E-utils.
\item Parsed with Biopython 1.83; computed GC directly, tallied features.
\item Ran independent CDS-name scan for canonical BGC markers.
\item Submitted CP018200.gb to antiSMASH v7.1.0 (web) --- 12 regions.
\item Ran antiSMASH v8.0.4 locally on uicgpu with \texttt{--cb-knownclusters --cb-subclusters --cc-mibig --clusterhmmer} --- 13 regions.
\item LLM-judge (\texttt{argo:gpt-5.2}, temp 0.1) on 21 claims + curated fact summary.
\end{enumerate}

\section{Named BGCs vs MIBiG (antiSMASH v8.0.4)}
\begin{longtable}{p{2.5cm}p{2.6cm}p{3.2cm}rp{1.6cm}r}
\toprule
Paper name & v8 region & MIBiG top hit & \#prot & \%id range & Score \\
\midrule
Difficidin & R2 (100--207\,kb) & BGC0000176 & 15 & 97--99 & 44658 \\
Fengycin & R5 (473--611\,kb) & BGC0001095 (+iturin, bacillomycin D, mycosubtilin) & 15 & 97--99 & 46028 \\
Bacillaene & R6 (676--786\,kb) & BGC0001089 & 14 & 97--99 & 46667 \\
Macrolactin & R7 (1005--1093\,kb) & BGC0000181 & 10 & 97--99 & 34862 \\
Class-II lanthipeptide & R8 (1259--1288\,kb) & \emph{none} & 0 & --- & --- \\
Surfactin & R11 ($\sim$2089--2154\,kb) & BGC0000433 & 21 & 95--100 & 29883 \\
Bacilysin & R12 (2776--2818\,kb) & BGC0001184 & 7 & 98--100 & 4632 \\
Bacillibactin & R13 (3354--3406\,kb) & BGC0000309 & 7 & 57--81 & 6049 \\
\bottomrule
\end{longtable}

\section{Verdict}
\textbf{PARTIAL $\rightarrow$ REPLICATED-leaning.} Every PGAP-derived genome-level number reproduces exactly. All seven named BGCs reconfirmed; six at $\geq$95\,\% pairwise identity, bacillibactin at 57--81\,\% (unambiguous cluster call, WS-8 uses the \textit{B.\ subtilis} dhbA-F variant, not the FZB42 variant). The class-II lanthipeptide is confirmed present and novel (no MIBiG hit). The BGC total count discrepancy (19 v3.0 vs 13 v8.0.4) is a tool-version artifact of well-documented antiSMASH stringency/merging changes. LC-MS and RNA-Seq claims are non-rerunnable because raw data were not deposited, but the genomic prerequisites for both are corroborated.

\section{GENUINE CRITIQUE}

\subsection{What genuinely replicated}
\begin{itemize}
\item Every quantitative genome-structure claim derived from the deposited GenBank flat file (bp length, topology, plasmid count, gene/CDS/pseudogene/tRNA/rRNA counts) reproduces exactly. GC\% direct computation matches the paper's \emph{abstract} value (46.5\,\%); the paper's Table 1 value (45.6\,\%) is almost certainly an internal typo/transposition.
\item All seven named BGC families (difficidin, fengycin, bacillaene, macrolactin, surfactin, bacilysin, bacillibactin) are unambiguously reconfirmed by fresh antiSMASH v8.0.4 + KnownClusterBlast against MIBiG, six of seven at $\geq$95\,\% per-gene identity.
\item The class-II lanthipeptide cluster is present, has the correct rule signature (LANC\_like + DUF4135), and returns no MIBiG hit --- consistent with the paper's own ``potential'' novelty framing.
\end{itemize}

\subsection{What genuinely did \emph{not} replicate}
\begin{itemize}
\item \textbf{Total BGC count.} Paper reports 19 (antiSMASH v3.0); this replication finds 13 (v8.0.4) or 12 (v7.1.0). The delta is not a paper error --- it is the direct consequence of antiSMASH's tightening of cluster-calling rules and expansion of the region-merging behaviour between v3 and v7/v8. It is nevertheless a real, non-trivial numerical discrepancy and should not be papered over: any downstream comparison against the paper's ``19'' number must carry the tool-version caveat explicitly.
\item \textbf{LC-MS metabolomics (C18).} The paper's core wet-lab claim --- that iturins and fengycins are the dominant antifungal principle --- cannot be numerically re-executed because the raw mass-spec data were never deposited (no MassIVE or MetaboLights accession is given). The best we can say is that the WS-8 genome \emph{has} the biosynthetic capacity to make both classes; whether they are actually the dominant species in the extract remains taken on trust.
\item \textbf{RNA-Seq expression (C19).} Similarly, the paper reports FPKM values for 27 core BGC genes but supplies no SRA/GEO accession. Genes are present; expression at the reported levels is not independently verifiable.
\item \textbf{Coverage claim ($\sim$311$\times$).} PacBio SMRT platform is confirmed via GenBank structured comment, but the raw reads (\texttt{fastq/bam}) are not linked to BioProject PRJNA354791, so depth cannot be recomputed.
\end{itemize}

\subsection{Systematic gaps and biases in this replication}
\begin{itemize}
\item \textbf{Species-identity not re-adjudicated.} We accepted the paper's ``\textit{B.\ amyloliquefaciens}'' species assignment without running independent digital DNA-DNA hybridization (dDDH) or average-nucleotide-identity (ANI) against the current \textit{B.\ amyloliquefaciens} / \textit{B.\ velezensis} / \textit{B.\ subtilis} type strains. Multiple 2015--2020 studies have shown that a large fraction of ``\textit{B.\ amyloliquefaciens} plantarum'' strains (including FZB42) have been reassigned to \textit{B.\ velezensis} under current TYGS/ANI thresholds. WS-8 was not re-classified in this replication --- an important open question (see \texttt{open\_questions.json} Q1).
\item \textbf{No independent re-assembly.} We used the deposited assembly rather than re-assembling from PacBio raw reads. Re-assembly is impossible in any case (raw reads absent from SRA), so any assembly-level artefact (misassembly, collapsed repeat, chimeric join at 3.93\,Mb) would remain undetected here.
\item \textbf{Bacillibactin identity gap not fully diagnosed.} The 57--81\,\% identity range for bacillibactin is called ``variant'' (\textit{B.\ subtilis} dhbA-F--like), but we did not perform a formal phylogenetic reconstruction of the \texttt{dhb} operon across \textit{Bacillus} spp.\ to confirm this.
\item \textbf{LLM-judge dependence.} The 83.6\,\% agreement score is generated by \texttt{argo:gpt-5.2} at temperature 0.1 --- a single model, single run, no bootstrap. The judgement is directional evidence, not a rigorous statistical estimate.
\end{itemize}

\subsection{What the replication tells us about the paper's real strength}
The paper's most defensible content is the genome + BGC inventory: those claims are the ones a third party can straightforwardly reconfirm from the deposited assembly, and they hold up. The paper's most vulnerable claims are the wet-lab ones that lack deposited raw data --- not because they are necessarily wrong, but because they are, six years later, unfalsifiable at the raw-data level. This asymmetry (strong genomics, unauditable wet-lab) is a typical pattern for 2020-era \textit{Bacillus} biocontrol papers and is a clean argument for stricter mandatory raw-data deposition (MassIVE, SRA) at submission.

\section{Provenance}
All artifacts under \texttt{\textasciitilde/Dropbox/REPLICATE-PROJECT/BVBRC-72-bamyloliquefaciens-ws8/}. Compute: CherryRd (orchestration, LLM judge via Argo proxy) + uicgpu (antiSMASH v8 local run, genome parsing). All endpoints FREE (NCBI, Europe PMC, Semantic Scholar, antiSMASH web, antiSMASH v8 local, Argo). No paid LLM API calls.

\end{document}
